%%%%%%%%%%%%%%%%Pakker og instillinger for dokumentet%%%%%%%%%%%%%%%%%

\documentclass[pdftex, 10pt, norsk, a4paper, twoside]{article} 
\usepackage[margin=1.3in]{geometry}							
\usepackage[norsk]{babel}%Legger til norske bokstaver		
\selectlanguage{norsk}%Får latex til å bruke norsk på steder der latex selv skriver ting. F.eks. Figur istedet for Figure	
\usepackage[T1]{fontenc} %Definerer hvilken font som skal brukes. For spesielt interreserte er dette en 8-bit encoding med fonter som bruker 256 glyfer							
\usepackage[utf8]{inputenc}%Tillater deg å skrive norske bokstaver direkte fra tastaturet, i steden for å bruke kommandoer.		

%%%%%%%% Figur-ting %%%%%%%%%%%%%
\usepackage{graphicx}% Nødvendig for å legge til bilder osv. 
\usepackage{wrapfig} %Gjør at du kan ha text wrapping rundt figurer
\usepackage{subcaption} %Gjør at du kan ha subfigurer
%%%%%%%%%%%%%%%%%%%%%%%%%%%%%%%

%%%%%%%%% SI-enheter %%%%%%%%%%%%
\usepackage{siunitx} % Lar deg bruke kommandoene \SI{tall}{\enhet} og \si{\enhet}. feks \si{\milli\gram\per\liter}. Dette er en veldig fin måte for å få konsistene enheter på. 
\sisetup{output-decimal-marker = {,}}%sikrer , og ikke . slik det skal være på norsk
\DeclareSIUnit{\bar}{bar} %Denne syntaksen kan du bruke for å lage nye SI-enheter
%%%%%%%%%%%%%%%%%%%%%%%%%%%%%%%%%

\usepackage[section]{placeins}
\usepackage{cancel}% for å kunne "stryke over" ting i ligninger feks \frac{\bcancel{b}}{\bcancel{b}} = 1
\usepackage{fmtcount}
\usepackage{booktabs}% lar deg bruke \toprule, \midrule og \bottomrule i tabeller for å lage pene tabeller.								
\usepackage{multirow}% Lar deg lage rader inni radene i en tabell.
\usepackage[pdftex,bookmarks,breaklinks]{hyperref} %legg til hidelinks hvis du ikke liker boksene rundt lenkene.
\usepackage{parskip}% Gjør at latex tolker paragrafer og ikke bare ignorerer "white space"
\usepackage{rotating} %Gir deg en rekke muligheter for å rotere ting
\usepackage{amsmath, amsfonts, amssymb}							
\usepackage[format=hang,font=footnotesize,labelfont=bf]{caption}% Gjør at captions under figurer og over tabeller blir pene.						
\usepackage[version=3]{mhchem} % Tillater deg å skrive kjemiske formler
\usepackage{epstopdf}% Omgjør .eps-filer til .pdf filer. Denne pakken er genial når man bruker chemdraw der man kan eksportere som eps. 
\usepackage{textcomp,gensymb}
\usepackage{float}% Tillater å bruke stor H i floats, dette kalles ofte for float of doom og er skjeldent anbefalt.
\usepackage{fancyhdr}% Lar deg lage den fine headeren på toppen av siden.

%%%%%%%%%%%%%%%%%%%%%%%%%%%%%%%%%%%%%%%%%%%%%%%%%%%%%%%%%%%%%
%%%%%%%%%% Vedlegg og kode %%%%%%%%%%%%%%%%%%%
\usepackage{minted} %pakke for å vise kodefiler
\usemintedstyle{perldoc} %penere stil
\usepackage{verbatim} % pakke for å vise .txt-filer
%%%%%%%%%%%%%%%%%%%%%%%%%%%%%%%%%%%%%%%%%%%%%%

%%%%%%%%%%% bibliografi %%%%%%%%%%%%%%%%%%%%%%%%%%%%%%%%%%%%%
\usepackage{csquotes} %nødvendig for at biblatex skal fungere med babel på norsk.
%sorting none gjør at referanselista kommer i siteringens rekkefølge. 
%natbib er en annen pakke, som jeg tror er mye utbredt i andres maler, men den er egentlig ganske utdatert, og funker ikke med norsk i det hele tatt, så jeg tok meg bryet med å finne ut hvordan biblatex funket slik at du ikke må:)
%\supercite i steden for bare \cite
\usepackage[backend=biber, style=ieee, sorting = none]{biblatex}
\addbibresource{referanser.bib}

%Denne endrer et par ord som babel har gjort skikkelig teite, slik at det ser litt penere ut.
\DefineBibliographyStrings{norsk}{
    urlseen = {Lest},
    url = {URL},
    andothers = {et al\adddot}
}
\let\oldcite\cite
\renewcommand{\cite}[1]{\textsuperscript{\oldcite{#1}}}
%Blokken over endrer \cite-kommandoen slik at den gir sitering i klammeparentes og opphøyd skrift, slik det skal være i vitenskapelige rapporter
%%%%%%%%%%%%%%%%%%%%%%%%%%%%%%%%%%%%%%%%%%%%%%%%%%%%%%%%%%%%%

%%%%%%%% Intern referering %%%%%%%%%%%%%%%%%%%%%%%%%%%%%%%%%%%
\usepackage{cleveref}%lar deg bruke kommando \Cref og \cref som legger inn ordet likning, tabell etc, med hhv stor og liten forbokstav.
\AddToHook{cmd/appendix/before}{\crefalias{section}{appendix}} % Gjør at sections i vedlegg kalles vedlegg av cref
\numberwithin{equation}{section}%Disse sørger for at tredje tabell i seksjon 2 blir (2.3) og ikke 15 eller noe sånt stygt.
\numberwithin{figure}{section}
\numberwithin{table}{section}
\crefname{appendix}{Vedlegg}{Vedlegg}
\crefname{section}{Seksjon}{Seksjonene}
%gjør figur til skjema, dette er strengt tatt et krav når du legger inn reaksjonsmekanismer. Kan fjernes om det er bilder du skal ha.
%\addto\captionsnorsk{\renewcommand{\figurename}{Skjema}}
%\crefname{figure}{Skjema}{Skjemaer}
%%%%%%%%%%%%%%%%%%%%%%%%%%%%%%%%%%%%%%%%%%%%%%%%%%%%%%%%%%%%%%%%%



\usepackage{soul}%lar deg bruke kommandoen \hl{tekst} for å gjøre midlertidig tekst rød
\sethlcolor{red}%velger at hl blir rød, og ikke feks lilla
\pagestyle{fancy}
\usepackage{diagbox} %lar deg ha skråstrek i tabeller øverst til venstre
\usepackage[toc, page]{appendix}
\usepackage{derivative} %gir deg korte kommandoer for derivasjon og partiellderivasjon
\usepackage{enumitem} 
\usepackage{chemfig} %kan skrive kjemiske formler
\usepackage{todonotes} %lar deg legge inn TODO-lapper i margen
\usepackage{xfrac} %lar deg ha skrå minibrøker i teksten
\usepackage{array} %lar deg lage nye kolonne-typer i tabeller
\newcolumntype{C}{>{$}c<{$}} %Stor C i kolonner i tabeller gir mattemodus på hele kolonna
%%%%%%%%%%%%%%%%%%%%%%%%%%%%%%%%%%%%%%%%%%%%%%%%%%%%%%%%%%%%%%%%%%%%%%

%%%%%%%%%%%%%%%%%%%%%%%%%%%%%%%%%%%%%%%%%%%%%%%%%%%%%%%%%%%%%%%%%%%%%%
%Setter noen instillinger for å gjøre ting pent :) 
\setlength\headheight{22.5pt}	
\floatstyle{plaintop}													
\restylefloat{table}													
\floatstyle{plain}														
\restylefloat{figure}													


%%%%%%%%%%%%%%%%%%%%%%%%%%%%%%%%%%%%%%%%%%%%%%%%%%%%%%%%%%%%%%%%%%%%%%
%Lager headeren
\lhead{Navn}			
\chead{}																
\rhead{Fagkode etc  \\ \today}	
\lfoot{}
\cfoot{\thepage}												
\rfoot{}
%%%%%%%%%%%%%%%%%%%%%%%%%%%%%%%%%%%%%%%%%%%%%%%%%%%%%%%%%%%%%%%%%%%%%%


\title{\textbf{\LaTeX-kurs}\\
\hspace{15pt}\\
\Large{Undertittel kan du ta her}}
\author{Navnet ditt her}
\date{\today}

%%%%%%%%%%%%%%%%Begynnelsen på dokumentet%%%%%%%%%%%%%%%%%
\begin{document}
% Tittel %
\maketitle
\thispagestyle{empty}
\newpage
% importerer fila som ligger i mappen Dokumenter og henter latexkoden derfra. Dette er for å ikke blande tekst og masse kode som er i dette dokumentet
\section{Matteting}

Dette er symbolet for delta $\Delta$
\begin{equation}
    4x + 2 = 5
    \label{eq:matteting test}
\end{equation}
\begin{equation}
    10 \cdot 5 = 50
    \label{eq:10x5}
\end{equation}
\begin{equation}
    \begin{split}
        a &= b + c \\
        &= c + d %&= gir = i midten
    \end{split}
    \label{eq:flere likninger}
\end{equation}
$10^{5x\cdot4^{54}}$ %krøllparanteser for ting som skal opp
$10_{4x\delta}$
$10^{\alpha}_{\epsilon}$

$\text{vanlig tekst} fisk$

$\alpha \hspace{5pt}\Gamma\hspace{5pt} \Sigma$
\begin{equation}
    (\frac{3x+5}{7x^{2}})
\end{equation}
\begin{equation}
    \left(\frac{3x+5}{7x^{2}}\right)
\end{equation}
\begin{equation}
    -\frac{\hbar^2}{2m} \frac{\partial^2\psi}{\partial x^2} =
    i \hbar \frac{\partial}{\partial t} \psi
    \label{eq:schrodinger} %ikke bruk norsk i labels
\end{equation}

\begin{equation}
    \int_{-\infty}^{\infty} \frac{1}{\sqrt{2\pi}\sigma}
    e^{-\frac{(X_i-\mu)^2}{2\sigma}}
    \label{eq:standard normalfordeling}
\end{equation}
\begin{equation*}
    2x^2 = 4x
    \label{eq:uten nummer}
\end{equation*}


\section{tabell}
\begin{table}[htb!]
    \centering
    \begin{tabular}{l c c | r}
    \toprule
        Prøve & Temperatur & Trykk & Kommentar  \\
        \# & [\si{\kilo\joule\per\mole}] & [\si{\pascal}] & tekst\\
        1 & 273 & $10^5$ & \\
    \midrule
        2 & & & \\ 
        3 & & & \\
    \bottomrule
    \end{tabular}
    \caption{Caption}
    \label{tab:my_label}
\end{table}
\section{Lister og generell skriving}
\section*{abstract} %stjerne gir ikkenummererte ting
\textbf{heihei} \textit{kursiv her} \textit{\textbf{begge deler}} Og vanlige kommandoer fungerer

\subsection{Punktliste}

\begin{itemize}
    \item Nummer én
    \item \LaTeX Kurs
    \item Chemikalen
    \begin{itemize}
        \item her er en underliste på denne
        \item Veldig gøy
    \end{itemize}
\end{itemize}

\subsection{Nummerert liste}
En liste med navn
\begin{enumerate}
    \item Nummer en
    \item Fire kalde pils
    \item I finland
    \begin{itemize}
        \item haha her kan man gjøre mye
    \end{itemize}
\end{enumerate}



\section{Kjemiting}
\ce{C2H5OH} \ce{CaCO3 (s) <=>[H2O] Ca^2+ + CO3^2+}

\section{SI-enheter}
\SI{1000}{\mol\per\cubic\centi\meter} %Stor gir deg tall og enheter

\si{\mol \per \cubic \meter\kilo\gram\squared} %liten si gir kun enhet

\si{\per \centi\meter}

\% gir deg prosenttegn likevel

\textcolor{red}{Her må du gjøre noe takk}
\input{Dokumenter/figurer.tex}

\newpage
%Referansene skal alltid komme på egen side
\printbibliography

\newpage 
%Ny side for vedlegg også
\appendix
\renewcommand{\thesection}{\Alph{section}} % endrer nummerering av seksjoner i vedlegg til bokstaver
\setcounter{page}{1} %setter sidetall på nytt
\pagenumbering{roman} %setter sidetall til romertall i vedlegg
\section{Vedlegg nr. 1}
\end{document} 